\documentclass[../../main.tex]{subfiles}
\begin{document}

{\bf [解]}

\begin{figure}[htbp]
  \centering
  \begin{tikzpicture}[scale=1.1, >=stealth, every node/.style={font=\small}]

    % --- 1. 格子と対角線（道路網） ---
    \foreach \x in {0,...,8} {
      \draw[gray!60] (\x, 0) -- (\x, 4);
    }
    \foreach \y in {0,...,4} {
      \draw[gray!60] (0, \y) -- (8, \y);
    }
    \foreach \x in {0,...,7} {
      \foreach \y in {0,...,3} {
        \draw[gray!50] (\x, \y) -- (\x+1, \y+1);
      }
    }

    % 外枠を強調
    \draw[thick] (0, 0) rectangle (8, 4);

    % --- 2. 始点・基準点 ---
    \fill (0, 4) circle (1.5pt);
    \node[left=2pt, font=\normalsize\bfseries] at (0, 4) {$A$};
    \fill (0, 0) circle (1.5pt);
    \node[left=2pt, font=\normalsize\bfseries] at (0, 0) {$B$};

    % --- 3. 点 P_j (到達境界の点) ---
    % 東側境界 (x = 8)
    \coordinate (P1) at (8, 0);
    \coordinate (P2) at (8, 1);
    \coordinate (P3) at (8, 2);
    \coordinate (P4) at (8, 3);
    \coordinate (P5) at (8, 4);

    % 北側境界 (y = 4)
    \coordinate (P6) at (7, 4);
    \coordinate (P7) at (6, 4);
    \coordinate (P8) at (5, 4);
    \coordinate (P9) at (4, 4);

    \foreach \j in {1,...,5} {
      \fill[blue!80!black] (P\j) circle (1.8pt);
    }
    \foreach \j in {6,...,9} {
      \fill[blue!80!black] (P\j) circle (1.8pt);
    }

    \node[below right] at (P1) {$P_1$};
    \node[right=2pt]   at (P2) {$P_2$};
    \node[right=2pt]   at (P3) {$P_3$};
    \node[right=2pt]   at (P4) {$P_4$};
    \node[above right] at (P5) {$P_5$};
    \node[above=2pt]   at (P6) {$P_6$};
    \node[above=2pt]   at (P7) {$P_7$};
    \node[above=2pt]   at (P8) {$P_8$};
    \node[above=2pt]   at (P9) {$P_9$};

    % --- 4. 点 Q_j (1 手手前の点) ---
    % x = 7 上 (東向き・北東向きの分岐点)
    \coordinate (Q1) at (7, 0);
    \coordinate (Q2) at (7, 1);
    \coordinate (Q3) at (7, 2);
    \coordinate (Q4) at (7, 3);

    % y = 3 上 (北東向きで北端へ抜ける直前の点)
    \coordinate (Q5) at (6, 3);
    \coordinate (Q6) at (5, 3);
    \coordinate (Q7) at (4, 3);
    \coordinate (Q8) at (3, 3);

    \foreach \j in {1,...,8} {
      \fill[red!80!black] (Q\j) circle (1.6pt);
    }

    % Q のラベル配置
    \node[below=2pt, red!80!black]       at (Q1) {$Q_1$};
    \node[below right=1pt, red!80!black] at (Q2) {$Q_2$};
    \node[below right=1pt, red!80!black] at (Q3) {$Q_3$};
    \node[above left=1pt, red!80!black]  at (Q4) {$Q_4$};
    \node[above left=1pt, red!80!black]  at (Q5) {$Q_5$};
    \node[above left=1pt, red!80!black]  at (Q6) {$Q_6$};
    \node[above left=1pt, red!80!black]  at (Q7) {$Q_7$};
    \node[above left=1pt, red!80!black]  at (Q8) {$Q_8$};

    % --- 5. 方位記号 ---
    \begin{scope}[shift={(1.0, -1.3)}]
      \draw[->, thick] (0, -0.6) -- (0, 0.6) node[above] {北};
      \draw[->, thick] (-0.6, 0) -- (0.6, 0) node[right] {東};
    \end{scope}

  \end{tikzpicture}
  \caption{道路網と各地点 $P_j$ および $Q_j$ の位置関係}
  \label{fig:grid_pj_qj}
\end{figure}

$P_j$ にくる確率 $a_j$，$Q_j$ にくる確率 $b_j$ とおくと，
題意の犯人の移動確率より，以下のような関係が成り立つ．

\begin{align}
  a_1 &= \left(\dfrac{1}{2}\right)^8, & a_5 &= \dfrac{1}{2} b_4 \\
  a_2 &= \dfrac{1}{2}(b_1 + b_2), & a_6 &= \dfrac{1}{2} b_5 \\
  a_3 &= \dfrac{1}{2}(b_2 + b_3), & a_7 &= \dfrac{1}{2} b_6 \\
  a_4 &= \dfrac{1}{2}(b_3 + b_4), & a_8 &= \dfrac{1}{2} b_7 \\
  & & a_9 &= \dfrac{1}{2} b_8
\end{align}

また，各$b_j$は何回北東方向に移動するかを考えることで以下のように場合の数・確率を計算できる．

\begin{table}[htb]
  \centering
  \caption{各点 $Q_j$ に到達する場合の数と確率 $b_j$}
  \begin{tabular}{c|c|c}
    $j$ & 場合の数 & 確率 $b_j$ \\ \hline
    $1$ & ${}_7\mathrm{C}_0 = 1$ & $(1/2)^7$ \\
    $2$ & ${}_7\mathrm{C}_1 = 7$ & $7(1/2)^7$ \\
    $3$ & ${}_7\mathrm{C}_2 = 21$ & $21(1/2)^7$ \\
    $4$ & ${}_7\mathrm{C}_3 = 35$ & $35(1/2)^7$ \\
    $5$ & ${}_6\mathrm{C}_3 = 20$ & $20(1/2)^6$ \\
    $6$ & ${}_5\mathrm{C}_2 = 10$ & $10(1/2)^5$ \\
    $7$ & ${}_4\mathrm{C}_1 = 4$ & $4(1/2)^4$ \\
    $8$ & $1$ & $(1/2)^3$
  \end{tabular}
\end{table}

従って，各 $a_j$ は以下のようになる ($\alpha = 1/2$)

\begin{table}[htb]
  \centering
  \caption{各点 $P_j$ に到達する確率 $a_j$ の一覧（$\alpha^8$ 倍した値）}
  \begin{tabular}{c|c || c|c}
    $j$ & $a_j \alpha^8$ & $j$ & $a_j \alpha^8$ \\ \hline
    $1$ & $1$ & $5$ & $35$ \\
    $2$ & $8$ & $6$ & $40$ \\
    $3$ & $28$ & $7$ & $40$ \\
    $4$ & $56$ & $8$ & $32$ \\
    & & $9$ & $16$
  \end{tabular}
\end{table}

従って，$a_j$ の最も大きい地点に警官を配置すればよく，
よって$P_4, P_6, P_7$ である．
この時，もとめる確率 $p$ は
\begin{align}
  p = \dfrac{56+40+40}{2^8}= \dfrac{136}{2^8} = \dfrac{17}{2^5} = 0.531 \fallingdotseq 0.53
\end{align}
である．

\end{document}
